% Original pages: 57-64
% Figures: none.
% Uncertainty log: none.

\noindent\textbf{25.} Let \(f:A\to B\), \(g:A\to C\) be ring homomorphisms and let \(h:A\to B\otimes_A C\) be defined by \(h(x)=f(x)\otimes g(x)\). Let \(X,Y,Z,T\) be the prime spectra of \(A,B,C,B\otimes_A C\) respectively. Then \(h^*(T)=f^*(Y)\cap g^*(Z)\).

\noindent[Let \(\mathfrak p\in X\), and let \(k=k(\mathfrak p)\) be the residue field at \(\mathfrak p\). By Exercise 21, the fiber \((h^*)^{-1}(\mathfrak p)\) is the spectrum of
\[
(B\otimes_A C)\otimes_A k\cong(B\otimes_A k)\otimes_k(C\otimes_A k).
\]
Hence \(\mathfrak p\in h^*(T)\Longleftrightarrow(B\otimes_A k)\otimes_k(C\otimes_A k)\neq0\Longleftrightarrow B\otimes_A k\neq0\) and \(C\otimes_A k\neq0\Longleftrightarrow\mathfrak p\in f^*(Y)\cap g^*(Z)\).]

\noindent\textbf{26.} Let \((B_\alpha,g_{\alpha\beta})\) be a direct system of rings and \(B\) the direct limit. For each \(\alpha\), let \(f_\alpha:A\to B_\alpha\) be a ring homomorphism such that \(g_{\alpha\beta}\circ f_\alpha=f_\beta\) whenever \(\alpha\leq\beta\) (i.e. the \(B_\alpha\) form a direct system of \(A\)-algebras). The \(f_\alpha\) induce \(f:A\to B\). Show that
\[
f^*(\operatorname{Spec}(B))=\bigcap_\alpha f_\alpha^*(\operatorname{Spec}(B_\alpha)).
\]

\noindent[Let \(\mathfrak p\in\operatorname{Spec}(A)\). Then \((f^*)^{-1}(\mathfrak p)\) is the spectrum of
\[
B\otimes_A k(\mathfrak p)\cong\varinjlim_\alpha(B_\alpha\otimes_A k(\mathfrak p))
\]
(since tensor products commute with direct limits: Chapter 2, Exercise 20). By Exercise 21 of Chapter 2 it follows that \((f^*)^{-1}(\mathfrak p)=\varnothing\) if and only if \(B_\alpha\otimes_A k(\mathfrak p)=0\) for some \(\alpha\), i.e., if and only if \((f_\alpha^*)^{-1}(\mathfrak p)=\varnothing\).]

\noindent\textbf{27.} i) Let \(f_\alpha:A\to B_\alpha\) be any family of \(A\)-algebras and let \(f:A\to B\) be their tensor product over \(A\) (Chapter 2, Exercise 23). Then
\[
f^*(\operatorname{Spec}(B))=\bigcap_\alpha f_\alpha^*(\operatorname{Spec}(B_\alpha)).
\]
[Use Examples 25 and 26.]

ii) Let \(f_\alpha:A\to B_\alpha\) be any finite family of \(A\)-algebras and let \(B=\prod_\alpha B_\alpha\). Define \(f:A\to B\) by \(f(x)=(f_\alpha(x))\). Then \(f^*(\operatorname{Spec}(B))=\bigcup_\alpha f_\alpha^*(\operatorname{Spec}(B_\alpha))\).

iii) Hence the subsets of \(X=\operatorname{Spec}(A)\) of the form \(f^*(\operatorname{Spec}(B))\), where \(f:A\to B\) is a ring homomorphism, satisfy the axioms for closed sets in a topological space. The associated topology is the \emph{constructible} topology on \(X\). It is finer than the Zariski topology (i.e., there are more open sets, or equivalently more closed sets).

iv) Let \(X_C\) denote the set \(X\) endowed with the constructible topology. Show that \(X_C\) is quasi-compact.

\noindent\textbf{28.} (Continuation of Exercise 27.)

i) For each \(g\in A\), the set \(X_g\) (Chapter 1, Exercise 17) is both open and closed in the constructible topology.

ii) Let \(C'\) denote the smallest topology on \(X\) for which the sets \(X_g\) are both open and closed, and let \(X_{C'}\) denote the set \(X\) endowed with this topology. Show that \(X_{C'}\) is Hausdorff.

iii) Deduce that the identity mapping \(X_C\to X_{C'}\) is a homeomorphism. Hence a subset \(E\) of \(X\) is of the form \(f^*(\operatorname{Spec}(B))\) for some \(f:A\to B\) if and only if it is closed in the topology \(C'\).

iv) The topological space \(X_C\) is compact, Hausdorff and totally disconnected.

\noindent\textbf{29.} Let \(f:A\to B\) be a ring homomorphism. Show that \(f^*:\operatorname{Spec}(B)\to\operatorname{Spec}(A)\) is a continuous \emph{closed} mapping (i.e., maps closed sets to closed sets) for the constructible topology.

\noindent\textbf{30.} Show that the Zariski topology and the constructible topology on \(\operatorname{Spec}(A)\) are the same if and only if \(A/\mathfrak N\) is absolutely flat (where \(\mathfrak N\) is the nilradical of \(A\)). [Use Exercise 11.]

\section{Primary Decomposition}

The decomposition of an ideal into primary ideals is a traditional pillar of ideal theory. It provides the algebraic foundation for decomposing an algebraic variety into its irreducible components---although it is only fair to point out that the algebraic picture is more complicated than na\"ive geometry would suggest. From another point of view primary decomposition provides a generalization of the factorization of an integer as a product of prime-powers. In the modern treatment, with its emphasis on localization, primary decomposition is no longer such a central tool in the theory. It is still, however, of interest in itself and in this chapter we establish the classical uniqueness theorems.

The prototypes of commutative rings are \(Z\) and the ring of polynomials \(k[x_1,\ldots,x_n]\) where \(k\) is a field; both these are unique factorization domains. This is not true of arbitrary commutative rings, even if they are integral domains (the classical example is the ring \(Z[\sqrt{-5}]\), in which the element 6 has two essentially distinct factorizations, \(2\cdot3\) and \((1+\sqrt{-5})(1-\sqrt{-5})\)). However, there is a generalized form of ``unique factorization'' of \emph{ideals} (not of elements) in a wide class of rings (the Noetherian rings).

A prime ideal in a ring \(A\) is in some sense a generalization of a prime number. The corresponding generalization of a power of a prime number is a primary ideal. An ideal \(\mathfrak q\) in a ring \(A\) is \emph{primary} if \(\mathfrak q\neq A\) and if
\[
xy\in\mathfrak q\quad\Longrightarrow\quad\text{either }x\in\mathfrak q\text{ or }y^n\in\mathfrak q\text{ for some }n>0.
\]
In other words,
\[
\mathfrak q\text{ is primary}\quad\Longleftrightarrow\quad A/\mathfrak q\neq0\text{ and every zero-divisor in }A/\mathfrak q\text{ is nilpotent}.
\]

Clearly every prime ideal is primary. Also the contraction of a primary ideal is primary, for if \(f:A\to B\) and if \(\mathfrak q\) is a primary ideal in \(B\), then \(A/\mathfrak q^c\) is isomorphic to a subring of \(B/\mathfrak q\).

\noindent\textbf{Proposition 4.1.\amtarget{4.1}} \emph{Let \(\mathfrak q\) be a primary ideal in a ring \(A\). Then \(r(\mathfrak q)\) is the smallest prime ideal containing \(\mathfrak q\).}

\noindent\emph{Proof.} By \amref{1.8} it is enough to show that \(\mathfrak p=r(\mathfrak q)\) is prime. Let \(xy\in r(\mathfrak q)\), then \((xy)^m\in\mathfrak q\) for some \(m>0\), and therefore either \(x^m\in\mathfrak q\) or \(y^{mn}\in\mathfrak q\) for some \(n>0\); i.e., either \(x\in r(\mathfrak q)\) or \(y\in r(\mathfrak q)\). \proofbox

If \(\mathfrak p=r(\mathfrak q)\), then \(\mathfrak q\) is said to be \(\mathfrak p\)-primary.

\noindent\textbf{Examples.} 1) The primary ideals in \(Z\) are \((0)\) and \((p^n)\), where \(p\) is prime. For these are the only ideals in \(Z\) with prime radical, and it is immediately checked that they are primary.

2) Let \(A=k[x,y]\), \(\mathfrak q=(x,y^2)\). Then \(A/\mathfrak q\cong k[y]/(y^2)\), in which the zero-divisors are all the multiples of \(y\), hence are nilpotent. Hence \(\mathfrak q\) is primary, and its radical \(\mathfrak p\) is \((x,y)\). We have \(\mathfrak p^2\subset\mathfrak q\subset\mathfrak p\) (strict inclusions), so that a primary ideal is not necessarily a prime-power.

3) Conversely, a prime power \(\mathfrak p^n\) is not necessarily primary, although its radical is the prime ideal \(\mathfrak p\). For example, let \(A=k[x,y,z]/(xy-z^2)\) and let \(\bar x,\bar y,\bar z\) denote the images of \(x,y,z\) respectively in \(A\). Then \(\mathfrak p=(\bar x,\bar z)\) is prime (since \(A/\mathfrak p\cong k[y]\), an integral domain); we have \(\bar x\bar y=\bar z^2\in\mathfrak p^2\) but \(\bar x\notin\mathfrak p^2\) and \(\bar y\notin r(\mathfrak p^2)=\mathfrak p\); hence \(\mathfrak p^2\) is not primary. However, there is the following result:

\noindent\textbf{Proposition 4.2.\amtarget{4.2}} \emph{If \(r(\mathfrak a)\) is maximal, then \(\mathfrak a\) is primary. In particular, the powers of a maximal ideal \(\mathfrak m\) are \(\mathfrak m\)-primary.}

\noindent\emph{Proof.} Let \(r(\mathfrak a)=\mathfrak m\). The image of \(\mathfrak m\) in \(A/\mathfrak a\) is the nilradical of \(A/\mathfrak a\), hence \(A/\mathfrak a\) has only one prime ideal, by \amref{1.8}. Hence every element of \(A/\mathfrak a\) is either a unit or nilpotent, and so every zero-divisor in \(A/\mathfrak a\) is nilpotent. \proofbox

We are going to study presentations of an ideal as an \emph{intersection of primary ideals}. First, a couple of lemmas:

\noindent\textbf{Lemma 4.3.\amtarget{4.3}} \emph{If \(\mathfrak q_i\) \((1\leq i\leq n)\) are \(\mathfrak p\)-primary, then \(\mathfrak q=\bigcap_{i=1}^n\mathfrak q_i\) is \(\mathfrak p\)-primary.}

\noindent\emph{Proof.} \(r(\mathfrak q)=r(\bigcap_{i=1}^n\mathfrak q_i)=\bigcap r(\mathfrak q_i)=\mathfrak p\). Let \(xy\in\mathfrak q\), \(y\notin\mathfrak q\). Then for some \(i\) we have \(xy\in\mathfrak q_i\) and \(y\notin\mathfrak q_i\), hence \(x\in\mathfrak p\), since \(\mathfrak q_i\) is primary. \proofbox

\noindent\textbf{Lemma 4.4.\amtarget{4.4}} \emph{Let \(\mathfrak q\) be a \(\mathfrak p\)-primary ideal, \(x\) an element of \(A\). Then}

\emph{i) if \(x\in\mathfrak q\) then \((\mathfrak q:x)=(1)\);}

\emph{ii) if \(x\notin\mathfrak q\) then \((\mathfrak q:x)\) is \(\mathfrak p\)-primary, and therefore \(r(\mathfrak q:x)=\mathfrak p\);}

\emph{iii) if \(x\notin\mathfrak p\) then \((\mathfrak q:x)=\mathfrak q\).}

\noindent\emph{Proof.} i) and iii) follow immediately from the definitions.

ii): if \(y\in(\mathfrak q:x)\) then \(xy\in\mathfrak q\), hence (as \(x\notin\mathfrak q\)) we have \(y\in\mathfrak p\). Hence \(\mathfrak q\subseteq(\mathfrak q:x)\subseteq\mathfrak p\); taking radicals, we get \(r(\mathfrak q:x)=\mathfrak p\). Let \(yz\in(\mathfrak q:x)\) with \(y\notin\mathfrak p\); then \(xyz\in\mathfrak q\), hence \(xz\in\mathfrak q\), hence \(z\in(\mathfrak q:x)\). \proofbox

A \emph{primary decomposition} of an ideal \(\mathfrak a\) in \(A\) is an expression of \(\mathfrak a\) as a finite intersection of primary ideals, say
\[
\mathfrak a=\bigcap_{i=1}^n\mathfrak q_i. \tag{1}
\]
(In general such a primary decomposition need not exist; in this chapter we shall restrict our attention to ideals which have a primary decomposition.) If moreover (i) the \(r(\mathfrak q_i)\) are all distinct, and (ii) we have \(\mathfrak q_i\not\supseteq\bigcap_{j\neq i}\mathfrak q_j\) \((1\leq i\leq n)\) the primary decomposition (1) is said to be \emph{minimal} (or irredundant, or reduced, or normal, \ldots). By \amref{4.3} we can achieve (i) and then we can omit any superfluous terms to achieve (ii); thus any primary decomposition can be reduced to a minimal one. We shall say that \(\mathfrak a\) is \emph{decomposable} if it has a primary decomposition.

\noindent\textbf{Theorem 4.5.\amtarget{4.5} (1st uniqueness theorem).} \emph{Let \(\mathfrak a\) be a decomposable ideal and let \(\mathfrak a=\bigcap_{i=1}^n\mathfrak q_i\) be a minimal primary decomposition of \(\mathfrak a\). Let \(\mathfrak p_i=r(\mathfrak q_i)\) \((1\leq i\leq n)\). Then the \(\mathfrak p_i\) are precisely the prime ideals which occur in the set of ideals \(r(\mathfrak a:x)\) \((x\in A)\), and hence are independent of the particular decomposition of \(\mathfrak a\).}

\noindent\emph{Proof.} For any \(x\in A\) we have \((\mathfrak a:x)=(\bigcap\mathfrak q_i:x)=\bigcap(\mathfrak q_i:x)\), hence \(r(\mathfrak a:x)=\bigcap_{i=1}^n r(\mathfrak q_i:x)=\bigcap_{x\notin\mathfrak q_j}\mathfrak p_j\) by \amref{4.4}. Suppose \(r(\mathfrak a:x)\) is prime; then by \amref{1.11} we have \(r(\mathfrak a:x)=\mathfrak p_j\) for some \(j\). Hence every prime ideal of the form \(r(\mathfrak a:x)\) is one of the \(\mathfrak p_j\). Conversely, for each \(i\) there exists \(x_i\notin\mathfrak q_i\), \(x_i\in\bigcap_{j\neq i}\mathfrak q_j\), since the decomposition is minimal; and we have \(r(\mathfrak a:x_i)=\mathfrak p_i\). \proofbox

\noindent\emph{Remarks.} 1) The above proof, coupled with the last part of \amref{4.4}, shows that for each \(i\) there exists \(x_i\) in \(A\) such that \((\mathfrak a:x_i)\) is \(\mathfrak p_i\)-primary.

2) Considering \(A/\mathfrak a\) as an \(A\)-module, \amref{4.5} is equivalent to saying that the \(\mathfrak p_i\) are precisely the prime ideals which occur as radicals of annihilators of elements of \(A/\mathfrak a\).

\noindent\textbf{Example.} Let \(\mathfrak a=(x^2,xy)\) in \(A=k[x,y]\). Then \(\mathfrak a=\mathfrak p_1\cap\mathfrak p_2^2\) where \(\mathfrak p_1=(x)\), \(\mathfrak p_2=(x,y)\). The ideal \(\mathfrak p_2^2\) is primary by \amref{4.2}. So the prime ideals are \(\mathfrak p_1,\mathfrak p_2\). In this example \(\mathfrak p_1\subset\mathfrak p_2\); we have \(r(\mathfrak a)=\mathfrak p_1\cap\mathfrak p_2=\mathfrak p_1\), but \(\mathfrak a\) is not a primary ideal.

The prime ideals \(\mathfrak p_i\) in \amref{4.5} are said to \emph{belong} to \(\mathfrak a\), or to be \emph{associated} with \(\mathfrak a\). The ideal \(\mathfrak a\) is primary if and only if it has only one associated prime ideal. The minimal elements of the set \(\{\mathfrak p_1,\ldots,\mathfrak p_n\}\) are called the \emph{minimal} or \emph{isolated} prime ideals belonging to \(\mathfrak a\). The others are called \emph{embedded} prime ideals. In the example above, \(\mathfrak p_2=(x,y)\) is embedded.

\noindent\textbf{Proposition 4.6.\amtarget{4.6}} \emph{Let \(\mathfrak a\) be a decomposable ideal. Then any prime ideal \(\mathfrak p\supseteq\mathfrak a\) contains a minimal prime ideal belonging to \(\mathfrak a\), and thus the minimal prime ideals of \(\mathfrak a\) are precisely the minimal elements in the set of all prime ideals containing \(\mathfrak a\).}

\noindent\emph{Proof.} If \(\mathfrak p\supseteq\mathfrak a=\bigcap_{i=1}^n\mathfrak q_i\), then \(\mathfrak p=r(\mathfrak p)\supseteq\bigcap r(\mathfrak q_i)=\bigcap\mathfrak p_i\). Hence by \amref{1.11} we have \(\mathfrak p\supseteq\mathfrak p_i\) for some \(i\); hence \(\mathfrak p\) contains a minimal prime ideal of \(\mathfrak a\). \proofbox

\noindent\emph{Remarks.} 1) The names \emph{isolated} and \emph{embedded} come from geometry. Thus if \(A=k[x_1,\ldots,x_n]\) where \(k\) is a field, the ideal \(\mathfrak a\) gives rise to a variety \(X\subseteq k^n\) (see Chapter 1, Exercise 25). The minimal primes \(\mathfrak p_i\) correspond to the irreducible components of \(X\), and the embedded primes correspond to subvarieties of these, i.e., varieties \emph{embedded} in the irreducible components. Thus in the example before \amref{4.6} the variety defined by \(\mathfrak a\) is the line \(x=0\), and the embedded ideal \(\mathfrak p_2=(x,y)\) corresponds to the origin \((0,0)\).

2) It is \emph{not} true that all the primary components are independent of the decomposition. For example \((x^2,xy)=(x)\cap(x,y)^2=(x)\cap(x^2,y)\) are two distinct minimal primary decompositions. However, there are some uniqueness properties: see \amref{4.10}.

\noindent\textbf{Proposition 4.7.\amtarget{4.7}} \emph{Let \(\mathfrak a\) be a decomposable ideal, let \(\mathfrak a=\bigcap_{i=1}^n\mathfrak q_i\) be a minimal primary decomposition, and let \(r(\mathfrak q_i)=\mathfrak p_i\). Then}
\[
\bigcup_{i=1}^n\mathfrak p_i=\{x\in A:(\mathfrak a:x)\neq\mathfrak a\}.
\]
\emph{In particular, if the zero ideal is decomposable, the set \(D\) of zero-divisors of \(A\) is the union of the prime ideals belonging to 0.}

\noindent\emph{Proof.} If \(\mathfrak a\) is decomposable, then 0 is decomposable in \(A/\mathfrak a\): namely \(0=\bigcap\bar{\mathfrak q}_i\) where \(\bar{\mathfrak q}_i\) is the image of \(\mathfrak q_i\) in \(A/\mathfrak a\), and is primary. Hence it is enough to prove the last statement of \amref{4.7}. By \amref{1.15} we have \(D=\bigcup_{x\neq0}r(0:x)\); from the proof of \amref{4.5}, we have \(r(0:x)=\bigcap_{x\notin\mathfrak q_j}\mathfrak p_j\subseteq\mathfrak p_j\) for some \(j\), hence \(D\subseteq\bigcup_{i=1}^n\mathfrak p_i\). But also from \amref{4.5} each \(\mathfrak p_i\) is of the form \(r(0:x)\) for some \(x\in A\), hence \(\bigcup\mathfrak p_i\subseteq D\). \proofbox

Thus (the zero ideal being decomposable)
\[
\begin{aligned}
D&=\text{set of zero-divisors}\\
 &=\bigcup\text{ of all prime ideals belonging to }0;\\
\mathfrak N&=\text{set of nilpotent elements}\\
 &=\bigcap\text{ of all minimal primes belonging to }0.
\end{aligned}
\]

Next we investigate the behavior of primary ideals under localization.

\noindent\textbf{Proposition 4.8.\amtarget{4.8}} \emph{Let \(S\) be a multiplicatively closed subset of \(A\), and let \(\mathfrak q\) be a \(\mathfrak p\)-primary ideal.}

\emph{i) If \(S\cap\mathfrak p\neq\varnothing\), then \(S^{-1}\mathfrak q=S^{-1}A\).}

\emph{ii) If \(S\cap\mathfrak p=\varnothing\), then \(S^{-1}\mathfrak q\) is \(S^{-1}\mathfrak p\)-primary and its contraction in \(A\) is \(\mathfrak q\). Hence primary ideals correspond to primary ideals in the correspondence \amref{3.11} between ideals in \(S^{-1}A\) and contracted ideals in \(A\).}

\noindent\emph{Proof.} i) If \(s\in S\cap\mathfrak p\), then \(s^n\in S\cap\mathfrak q\) for some \(n>0\); hence \(S^{-1}\mathfrak q\) contains \(s^n/1\), which is a unit in \(S^{-1}A\).

ii) If \(S\cap\mathfrak p=\varnothing\), then \(s\in S\) and \(as\in\mathfrak q\) imply \(a\in\mathfrak q\), hence \(\mathfrak q^{ec}=\mathfrak q\) by \amref{3.11}. Also from \amref{3.11} we have \(r(\mathfrak q^e)=r(S^{-1}\mathfrak q)=S^{-1}r(\mathfrak q)=S^{-1}\mathfrak p\). The verification that \(S^{-1}\mathfrak q\) is primary is straightforward. Finally, the contraction of a primary ideal is primary. \proofbox

For any ideal \(\mathfrak a\) and any multiplicatively closed subset \(S\) in \(A\), the contraction in \(A\) of the ideal \(S^{-1}\mathfrak a\) is denoted by \(S(\mathfrak a)\).

\noindent\textbf{Proposition 4.9.\amtarget{4.9}} \emph{Let \(S\) be a multiplicatively closed subset of \(A\) and let \(\mathfrak a\) be a decomposable ideal. Let \(\mathfrak a=\bigcap_{i=1}^n\mathfrak q_i\) be a minimal primary decomposition of \(\mathfrak a\). Let \(\mathfrak p_i=r(\mathfrak q_i)\) and suppose the \(\mathfrak q_i\) numbered so that \(S\) meets \(\mathfrak p_{m+1},\ldots,\mathfrak p_n\) but not \(\mathfrak p_1,\ldots,\mathfrak p_m\). Then}
\[
S^{-1}\mathfrak a=\bigcap_{i=1}^mS^{-1}\mathfrak q_i,
\qquad
S(\mathfrak a)=\bigcap_{i=1}^m\mathfrak q_i,
\]
\emph{and these are minimal primary decompositions.}

\noindent\emph{Proof.} \(S^{-1}\mathfrak a=\bigcap_{i=1}^nS^{-1}\mathfrak q_i\) by \amref{3.11} \(=\bigcap_{i=1}^mS^{-1}\mathfrak q_i\) by \amref{4.8}, and \(S^{-1}\mathfrak q_i\) is \(S^{-1}\mathfrak p_i\)-primary for \(i=1,\ldots,m\). Since the \(\mathfrak p_i\) are distinct, so are the \(S^{-1}\mathfrak p_i\) \((1\leq i\leq m)\), hence we have a minimal primary decomposition. Contracting both sides, we get
\[
S(\mathfrak a)=(S^{-1}\mathfrak a)^c=\bigcap_{i=1}^m(S^{-1}\mathfrak q_i)^c=\bigcap_{i=1}^m\mathfrak q_i
\]
by \amref{4.8} again. \proofbox

A set \(\Sigma\) of prime ideals belonging to \(\mathfrak a\) is said to be \emph{isolated} if it satisfies the following condition: if \(\mathfrak p'\) is a prime ideal belonging to \(\mathfrak a\) and \(\mathfrak p'\subseteq\mathfrak p\) for some \(\mathfrak p\in\Sigma\), then \(\mathfrak p'\in\Sigma\).

Let \(\Sigma\) be an isolated set of prime ideals belonging to \(\mathfrak a\), and let \(S=A-\bigcup_{\mathfrak p\in\Sigma}\mathfrak p\). Then \(S\) is multiplicatively closed and, for any prime ideal \(\mathfrak p'\) belonging to \(\mathfrak a\), we have
\[
\mathfrak p'\in\Sigma\quad\Longrightarrow\quad\mathfrak p'\cap S=\varnothing;
\]
\[
\mathfrak p'\notin\Sigma\quad\Longrightarrow\quad\mathfrak p'\not\subseteq\bigcup_{\mathfrak p\in\Sigma}\mathfrak p\quad\text{(by \amref{1.11})}\quad\Longrightarrow\quad\mathfrak p'\cap S\neq\varnothing.
\]
Hence, from \amref{4.9}, we deduce

\noindent\textbf{Theorem 4.10.\amtarget{4.10} (2nd uniqueness theorem).} \emph{Let \(\mathfrak a\) be a decomposable ideal, let \(\mathfrak a=\bigcap_{i=1}^n\mathfrak q_i\) be a minimal primary decomposition of \(\mathfrak a\), and let \(\{\mathfrak p_{i_1},\ldots,\mathfrak p_{i_m}\}\) be an isolated set of prime ideals of \(\mathfrak a\). Then \(\mathfrak q_{i_1}\cap\cdots\cap\mathfrak q_{i_m}\) is independent of the decomposition.}

In particular:

\noindent\textbf{Corollary 4.11.\amtarget{4.11}} \emph{The isolated primary components (i.e., the primary components \(\mathfrak q_i\) corresponding to minimal prime ideals \(\mathfrak p_i\)) are uniquely determined by \(\mathfrak a\).}

\noindent\emph{Proof of \amref{4.10}.} We have \(\mathfrak q_{i_1}\cap\cdots\cap\mathfrak q_{i_m}=S(\mathfrak a)\) where \(S=A-(\mathfrak p_{i_1}\cup\cdots\cup\mathfrak p_{i_m})\), hence depends only on \(\mathfrak a\) (since the \(\mathfrak p_i\) depend only on \(\mathfrak a\)). \proofbox

\noindent\emph{Remark.} On the other hand, the embedded primary components are in general not uniquely determined by \(\mathfrak a\). If \(A\) is a Noetherian ring, there are in fact infinitely many choices for each embedded component (see Chapter 8, Exercise 1).

\subsection{Exercises}

\noindent\textbf{1.} If an ideal \(\mathfrak a\) has a primary decomposition, then \(\operatorname{Spec}(A/\mathfrak a)\) has only finitely many irreducible components.

\noindent\amstar\textbf{2.} If \(\mathfrak a=r(\mathfrak a)\), then \(\mathfrak a\) has no embedded prime ideals.

\noindent\textbf{3.} If \(A\) is absolutely flat, every primary ideal is maximal.

\noindent\textbf{4.} In the polynomial ring \(Z[t]\), the ideal \(\mathfrak m=(2,t)\) is maximal and the ideal \(\mathfrak q=(4,t)\) is \(\mathfrak m\)-primary, but is not a power of \(\mathfrak m\).

\noindent\textbf{5.} In the polynomial ring \(K[x,y,z]\) where \(K\) is a field and \(x,y,z\) are independent indeterminates, let \(\mathfrak p_1=(x,y)\), \(\mathfrak p_2=(x,z)\), \(\mathfrak m=(x,y,z)\); \(\mathfrak p_1\) and \(\mathfrak p_2\) are prime, and \(\mathfrak m\) is maximal. Let \(\mathfrak a=\mathfrak p_1\mathfrak p_2\). Show that \(\mathfrak a=\mathfrak p_1\cap\mathfrak p_2\cap\mathfrak m^2\) is a reduced primary decomposition of \(\mathfrak a\). Which components are isolated and which are embedded?

\noindent\textbf{6.} Let \(X\) be an infinite compact Hausdorff space, \(C(X)\) the ring of real-valued continuous functions on \(X\) (Chapter 1, Exercise 26). Is the zero ideal decomposable in this ring?

\noindent\textbf{7.} Let \(A\) be a ring and let \(A[x]\) denote the ring of polynomials in one indeterminate over \(A\). For each ideal \(\mathfrak a\) of \(A\), let \(\mathfrak a[x]\) denote the set of all polynomials in \(A[x]\) with coefficients in \(\mathfrak a\).

i) \(\mathfrak a[x]\) is the extension of \(\mathfrak a\) to \(A[x]\).

ii) If \(\mathfrak p\) is a prime ideal in \(A\), then \(\mathfrak p[x]\) is a prime ideal in \(A[x]\).

iii) If \(\mathfrak q\) is a \(\mathfrak p\)-primary ideal in \(A\), then \(\mathfrak q[x]\) is a \(\mathfrak p[x]\)-primary ideal in \(A[x]\). [Use Chapter 1, Exercise 2.]

iv) If \(\mathfrak a=\bigcap_{i=1}^n\mathfrak q_i\) is a minimal primary decomposition in \(A\), then \(\mathfrak a[x]=\bigcap_{i=1}^n\mathfrak q_i[x]\) is a minimal primary decomposition in \(A[x]\).

v) If \(\mathfrak p\) is a minimal prime ideal of \(\mathfrak a\), then \(\mathfrak p[x]\) is a minimal prime ideal of \(\mathfrak a[x]\).

\noindent\textbf{8.} Let \(k\) be a field. Show that in the polynomial ring \(k[x_1,\ldots,x_n]\) the ideals \(\mathfrak p_i=(x_1,\ldots,x_i)\) \((1\leq i\leq n)\) are prime and all their powers are primary. [Use Exercise 7.]

\noindent\textbf{9.} In a ring \(A\), let \(D(A)\) denote the set of prime ideals \(\mathfrak p\) which satisfy the following condition: there exists \(a\in A\) such that \(\mathfrak p\) is minimal in the set of prime ideals containing \((0:a)\). Show that \(x\in A\) is a zero divisor \(\Longleftrightarrow x\in\mathfrak p\) for some \(\mathfrak p\in D(A)\).

Let \(S\) be a multiplicatively closed subset of \(A\), and identify \(\operatorname{Spec}(S^{-1}A)\) with its image in \(\operatorname{Spec}(A)\) (Chapter 3, Exercise 21). Show that
\[
D(S^{-1}A)=D(A)\cap\operatorname{Spec}(S^{-1}A).
\]
If the zero ideal has a primary decomposition, show that \(D(A)\) is the set of associated prime ideals of 0.

\noindent\textbf{10.} For any prime ideal \(\mathfrak p\) in a ring \(A\), let \(S_\mathfrak p(0)\) denote the kernel of the homomorphism \(A\to A_\mathfrak p\). Prove that

i) \(S_\mathfrak p(0)\subseteq\mathfrak p\).

ii) \(r(S_\mathfrak p(0))=\mathfrak p\Longleftrightarrow\mathfrak p\) is a minimal prime ideal of \(A\).

iii) If \(\mathfrak p\supseteq\mathfrak p'\), then \(S_\mathfrak p(0)\subseteq S_{\mathfrak p'}(0)\).

iv) \(\bigcap_{\mathfrak p\in D(A)}S_\mathfrak p(0)=0\), where \(D(A)\) is defined in Exercise 9.
